\chapter[\hspace{0.6em}时变处理]{时变处理}

% 本文件内的局部辅助：带垂直偏移的边注，用于错开相邻批注
\providecommand{\mnotev}[2]{\marginnote{\textcolor{black!85}{#2}}[#1]}

到目前为止，本书讨论的都是不随时间变化的固定处理。然而，许多因果问题所涉及的处理会随时间变化。例如，我们可能想估计医疗处理、生活习惯、就业状况、婚姻状况、职业暴露等的因果效应。由于这些处理对同一个人在不同时间可能取不同取值，我们称之为时变处理。

在本书第一、第二部分，我们把注意力限制在时不变处理上，这有助于引入基本概念与方法。现在该考虑更现实的因果问题了——它们涉及在时间中逐步展开的假想干预之间的对比。第三部分把第一、第二部分的内容推广到时变处理。本章介绍时变处理因果推断的一些关键术语与概念。尽管我们已尽力简化这些概念（不信的话，去翻翻因果推断文献），本章仍是全书技术性最强的章节之一。遗憾的是，再作简化就会损失太多严谨性。不过，既然你读到了这里，你就有资格读懂本章。

\section{时变处理的因果效应}

考虑随访起点（时点零）上的一个时不变处理变量 $A$（1：受处理，0：未处理），以及 60 个月后测量的结局变量 $Y$。我们此前把 $A$ 对结局 $Y$ 的平均因果效应定义为：处理之下的平均反事实结局 $\E[Y^{a=1}]$ 与未处理之下的平均反事实结局 $\E[Y^{a=0}]$ 之间的对比，即 $\E[Y^{a=1}]-\E[Y^{a=0}]$。由于每个人的处理状态都在单一时刻（时点零）确定，平均因果效应无需提及处理发生的时间。相反，涉及时变处理的因果对比则必须把时间明确纳入。

为看清这一点，考虑一个二值时变处理 $A_k$，它可能在随访的每个月 $k$ 变化，其中 $k=0,1,2,\ldots,K$，$K=59$。例如，在一项对感染人类免疫缺陷病毒（HIV）者的 5 年随访研究中，若个体在第 $k$ 个月接受抗逆转录病毒治疗，则 $A_k$ 取 1，否则取 0。在研究开始时（时点 0）之前没有人接受治疗，即对所有个体都有 $A_{-1}=0$。

\mnotev{-0.4em}{为简单起见，我们暂时假定此期间没有人失访，并假定所有变量都被完美测量。}

我们用上横线表示处理史，即 $\bar{A}_k=(A_0,A_1,\ldots,A_k)$ 是从时点 0 到时点 $k$ 的处理史。当我们指涉贯穿到 $K$ 的整个处理史时，常把 $\bar{A}_K$ 写作不带下标的 $\bar{A}$。在我们的 HIV 研究中，整个随访期间持续接受治疗的个体，其处理史为 $\bar{A}=(A_0=1,A_1=1,\ldots,A_{59}=1)=(1,1,\ldots,1)$，或 $\bar{A}=\bar{1}$。类似地，随访期间从未接受治疗的个体，处理史为 $\bar{A}=(0,0,\ldots,0)=\bar{0}$。多数个体只在随访的部分时间接受治疗，因而其处理史是既有 1 又有 0 的中间处理史——我们无法像 $\bar{1}$ 和 $\bar{0}$ 那样紧凑地表示它们。

\mnotev{-0.4em}{为与许多已发表论文保持一致，我们对时间采用从 0 开始的编号。也就是说，可能接受治疗的第一个时点是 $k=0$ 而不是 $k=1$。}

设 $Y$ 测量随访结束时（时点 $K+1=60$）的健康状况——$Y$ 取值越大表示健康越好。我们想估计时变处理 $\bar{A}$ 对结局 $Y$ 的平均因果效应。但我们不能再把时变处理的平均因果效应定义为单一时点 $k$ 上的对比，因为对比 $\E[Y^{a_k=1}]-\E[Y^{a_k=0}]$ 量度的是单一时点 $k$ 上处理 $A_k$ 的效应，而不是 0 到 59 之间所有时点上的时变处理 $A_k$ 的效应。事实上，我们必须把平均因果效应定义为两种处理策略之下的反事实平均结局之间的对比，而这两种策略都涉及从随访开始（$k=0$）到结束（$k=K$）所有时点上的处理。由此，时变处理的平均因果效应不是唯一确定的。下一节我们描述时变处理平均因果效应的许多可能定义。

\mnotev{-0.4em}{为简单起见，我们的例子考虑在固定时点测量的结局。不过，本章讨论的概念也适用于时变结局与失效时间结局（见技术注记 21.8）。}

\section{处理策略}

处理策略——也称计划、政策、规程或方案（regime）——是在随访的每个时点 $k$ 分配处理的规则。例如，两种处理策略是随访期间的“始终治疗”与“从不治疗”。“始终治疗”策略表示为 $\bar{a}=(1,1,\ldots,1)=\bar{1}$，“从不治疗”策略表示为 $\bar{a}=(0,0,\ldots,0)=\bar{0}$。于是我们可以把 $\bar{A}$ 对结局 $Y$ 的平均因果效应定义为“始终治疗”策略下的平均反事实结局 $\E[Y^{\bar{a}=\bar{1}}]$ 与“从不治疗”策略下的平均反事实结局 $\E[Y^{\bar{a}=\bar{0}}]$ 之间的对比，即
\[
\E\bigl[Y^{\bar{a}=\bar{1}}\bigr]-\E\bigl[Y^{\bar{a}=\bar{0}}\bigr].
\]

\mnotev{-0.4em}{比较处理策略的一般反事实理论，最早由 Robins (1986, 1987, 1997a) 提出。}

但对时变处理 $\bar{A}$ 还有许多其他可能的因果效应，每一个都由两种特定处理策略下的结局对比来定义。例如，我们可能对由对比 $\E[Y^{\bar{a}}]-\E[Y^{\bar{a}'}]$ 定义的平均因果效应感兴趣，它比较“每隔一个月治疗一次”的策略 $\bar{a}=(1,0,1,0,\ldots)$ 与“除第一个月外每个月都治疗”的策略 $\bar{a}'=(0,1,1,1,\ldots)$。可能的对比数量极大：对二值 $a_k$，取值 $(a_0,a_1,\ldots,a_K)$ 有 $2^K$ 种可能组合，所以我们至少可以定义 $2^K$ 种处理策略。事实上，正如我们下面要说明的，这 $2^K$ 种策略并未穷尽所有可能的处理策略。

\mnotev{-0.2em}{记住我们用小写字母表示随机变量的可能实现：$a_k$ 是处理 $A_k$ 的一个实现。}

为在 HIV 例子中定义更多处理策略，考虑时变协变量 $L_k$，它表示在第 $k$ 个月对所有人测量的 CD4 细胞计数（单位：个/$\mu$L）。当 CD4 细胞计数低（提示预后不良）时 $L_k$ 取 1，否则取 0。在时点零，所有人的 CD4 细胞计数都高，$L_0=0$。于是我们可以考虑策略“当 $L_k=0$ 时不治疗，当 $L_k=1$ 时开始治疗，此后持续治疗”。这一处理策略不同于上一段考虑的那些，因为我们无法用规则 $\bar{a}=(a_0,a_1,\ldots,a_K)$ 来表示它——在那种规则下，所有个体在 $k=0$ 时都接受同样的处理 $a_0$，在 $k=1$ 时都接受 $a_1$，如此等等。而在这里，每个时点上有些个体会受治疗、另一些不受治疗，取决于其不断变化的 $L_k$ 的取值。这是动态处理策略的一个例子：在时点 $k$ 的处理 $a_k$ 取决于个体时变协变量 $\bar{L}_k$ 的演化。处理不依赖协变量的策略 $\bar{a}$ 称为非动态策略或静态策略。正式定义见精要 19.1。

\begin{finepoint}[确定性与随机处理策略]
确定性动态处理策略是一个规则 $g=\bigl[g_0(\bar{a}_{-1},\bar{l}_0),\ldots,g_K(\bar{a}_{K-1},\bar{l}_K)\bigr]$，其中 $g_k(\bar{a}_{k-1},\bar{l}_k)$ 指定给具有过去史 $(\bar{a}_{k-1},\bar{l}_k)$ 的个体在 $k$ 时分配的处理。HIV 研究中的一个例子：若个体的 CD4 细胞计数（$\bar{l}_k$ 的函数）在 $k$ 时或之前一直很低，则 $g_k(\bar{a}_{k-1},\bar{l}_k)$ 为 1，否则为 0。确定性静态处理策略是一个规则 $g=\bigl[g_0(\bar{a}_{-1}),\ldots,g_K(\bar{a}_{K-1})\bigr]$，其中 $g_k(\bar{a}_{k-1})$ 不依赖 $\bar{l}_k$。我们常把 $g_k(\bar{a}_{k-1},\bar{l}_k)$ 简写为 $g_k$。

我们想比较的多数静态与动态策略都是确定性处理策略：它们在每个时点给每个个体分配一个特定的处理取值（0 或 1）。更一般地，我们可以考虑随机处理策略：它不分配特定的处理取值，而是分配接受某处理取值的概率。随机处理策略可以是静态的（例如“每个月独立地以概率 0.3 治疗、以概率 0.7 不治疗”），也可以是动态的（例如“每个月独立地以概率 0.3 治疗 CD4 细胞计数低的个体，但不治疗 CD4 细胞计数高的个体”）。

我们把使平均反事实结局 $\E[Y^g]$ 最大（当结局取值越大越好时）的策略 $g$ 称为最优处理策略。对药物治疗而言，最优策略几乎总是动态的，因为出现毒性时必须停药。此外，任何随机策略都不可能优于最优确定性策略。然而，随机策略（即随机化试验）在科学上仍是必要的，因为在试验之前，无人知道哪个确定性策略最优。关于处理策略的分类，见 Young et al. (2014)。在正文中，除另有说明外，字母 $g$ 只指确定性处理策略。
\end{finepoint}

时变处理的因果推断涉及两种或更多处理策略之下反事实结局的对比。只有当所关心的处理策略被明确指定时，时变处理的平均因果效应才是定义良好的。在 HIV 例子中，我们可以基于差值 $\E[Y^{\bar{a}}]-\E[Y^{\bar{a}'}]$ 定义平均因果效应，它对比策略 $\bar{a}$（比如“始终治疗”）与策略 $\bar{a}'$（比如“从不治疗”）；也可以基于差值 $\E[Y^{\bar{a}}]-\E[Y^{g}]$ 定义，它对比策略 $\bar{a}$（“始终治疗”）与策略 $g$（比如“仅在 CD4 细胞计数变低后才治疗”）。注意我们常用 $g$ 表示任何策略——静态的或动态的。当我们用它表示静态策略时，有时写作 $Y^{g=\bar{a}}$ 而不是 $Y^g$ 或 $Y^{\bar{a}}$。

也就是说，时变处理的因果效应没有唯一的定义。即便每个时点 $k$ 上只有两种处理选项——治疗或不治疗——我们仍可定义与处理策略对一样多的因果效应。下一节我们描述一种研究设计，在它之下所有这些因果效应都能被有效估计：序贯随机化实验。

\section{序贯随机化实验}

图 19.1、19.2 与 19.3 中的因果图概括了时变处理研究中可能出现的三种情形。在三个图中，$A_k$ 表示时变处理，$L_k$ 表示已测量变量集合，$Y$ 是结局，$U_k$ 是 $k$ 时至少另外两个变量的共同原因（即不可测量变量集合）。由于协变量 $U_k$ 未被测量，其取值未知，因而无法用于分析。在我们的 HIV 研究中，时变 CD4 细胞计数 $L_k$ 是真实但未测量的免疫系统慢性损伤 $U_k$ 的结果。个体的免疫损伤 $U_k$ 越重，其 CD4 细胞计数 $L_k$ 与健康状况 $Y$ 就越低。为简单起见，因果图只画出前两个随访时点 $k=0$ 与 $k=1$，并假定所有参与者都依从所分配的处理。

图 19.1 的因果图中，既没有从已测量协变量 $\bar{L}_k$ 指向处理 $A_k$ 的箭头，也没有从不可测量协变量 $\bar{U}_k$ 指向处理 $A_k$ 的箭头。图 19.2 的因果图有从已测量协变量 $\bar{L}_k$ 指向处理 $A_k$ 的箭头，但没有从不可测量协变量 $\bar{U}_k$ 指向处理 $A_k$ 的箭头。图 19.3 的因果图既有从已测量协变量 $\bar{L}_k$、也有从不可测量协变量 $\bar{U}_k$ 指向处理 $A_k$ 的箭头。

\begin{figure}[htbp]
\centering
\begin{tikzpicture}[causal,every node/.style={transform shape}]
  \node (L0) at (0,1.35) {$L_0$};
  \node (A0) at (2.1,1.35) {$A_0$};
  \node (L1) at (4.3,1.35) {$L_1$};
  \node (A1) at (6.5,1.35) {$A_1$};
  \node (Y)  at (8.7,1.35) {$Y$};
  \node (U0) at (0,0) {$U_0$};
  \node (U1) at (4.3,0) {$U_1$};
  \draw[ed] (U0) -- (L0);
  \draw[ed] (U0) -- (U1);
  \draw[ed] (U1) -- (L1);
  \draw[ed] (A0) -- (L1);
  \draw[ed] (A1) -- (Y);
  \draw[ed] (A0) -- (U1);
  \draw[ed] (U0) .. controls (2.2,0.10) and (3.3,0.85) .. (L1);
  \draw[ed] (L0) .. controls (1.2,2.10) and (3.1,2.10) .. (L1);
  \draw[ed] (A0) .. controls (3.1,2.50) and (5.5,2.50) .. (A1);
  \draw[ed] (L1) .. controls (5.4,2.00) and (7.3,2.00) .. (Y);
  \draw[ed] (A0) .. controls (4.3,3.05) and (6.6,3.05) .. (Y);
  \draw[ed] (L0) .. controls (2.8,3.60) and (6.2,3.60) .. (Y);
  \draw[ed] (U0) .. controls (2.8,-1.00) and (6.2,-1.00) .. (Y);
  \draw[ed] (U1) .. controls (5.6,-0.55) and (7.4,-0.55) .. (Y);
\end{tikzpicture}
\caption{}
\label{fig:19.1}
\end{figure}

图 19.1 可以表示一项随机化实验：每个时点 $k$ 的处理 $A_k$ 都按只依赖既往处理史的概率随机分配（为简单起见，我们通篇假定完美依从）。若例如个体在第 $k$ 个月的处理取值这样随机分配——上月未接受治疗者（$A_{k-1}=0$）以概率 0.5 受治疗，上月已接受治疗者（$A_{k-1}=1$）以概率 1 受治疗——则我们的 HIV 研究就由图 19.1 表示。当我们关心静态处理策略的对比时，图 19.1 是“不存在由已测量或不可测量变量引起的混杂”这一条件在时变处理情形的恰当推广。在该因果图下，若所有人都遵循静态处理策略 $\bar{a}$，则反事实结局均值 $\E[Y^{\bar{a}}]$ 就是遵循策略 $\bar{a}$ 者的平均结局 $\E[Y\given \bar{A}=\bar{a}]$。（有趣的是，对动态策略这并不成立。对依赖变量 $L$ 的动态策略 $g$，只有当处理 $A_k$ 依赖 $L_k$ 的所有时点 $k$ 上接受处理 $A_k=1$ 的概率恰为 0.5 时，反事实均值 $\E[Y^g]$ 才等于遵循策略 $g$ 者的平均结局。否则，识别 $\E[Y^g]$ 需要在图 19.1 或图 19.2 之下对 $L$、$A$ 与 $Y$ 的数据应用 g 方法。）

\begin{figure}[htbp]
\centering
\begin{tikzpicture}[causal,every node/.style={transform shape}]
  \node (L0) at (0,1.35) {$L_0$};
  \node (A0) at (2.1,1.35) {$A_0$};
  \node (L1) at (4.3,1.35) {$L_1$};
  \node (A1) at (6.5,1.35) {$A_1$};
  \node (Y)  at (8.7,1.35) {$Y$};
  \node (U0) at (0,0) {$U_0$};
  \node (U1) at (4.3,0) {$U_1$};
  \draw[ed] (U0) -- (L0);
  \draw[ed] (U0) -- (U1);
  \draw[ed] (U1) -- (L1);
  \draw[ed] (L0) -- (A0);
  \draw[ed] (A0) -- (L1);
  \draw[ed] (L1) -- (A1);
  \draw[ed] (A1) -- (Y);
  \draw[ed] (A0) -- (U1);
  \draw[ed] (U0) .. controls (2.2,0.10) and (3.3,0.85) .. (L1);
  \draw[ed] (L0) .. controls (1.2,2.10) and (3.1,2.10) .. (L1);
  \draw[ed] (L0) .. controls (4.6,2.95) and (5.9,2.35) .. (A1);
  \draw[ed] (A0) .. controls (3.1,2.50) and (5.5,2.50) .. (A1);
  \draw[ed] (L1) .. controls (5.4,2.00) and (7.3,2.00) .. (Y);
  \draw[ed] (A0) .. controls (4.5,3.35) and (6.8,3.35) .. (Y);
  \draw[ed] (L0) .. controls (2.9,3.85) and (6.3,3.85) .. (Y);
  \draw[ed] (U0) .. controls (2.8,-1.00) and (6.2,-1.00) .. (Y);
  \draw[ed] (U1) .. controls (5.6,-0.55) and (7.4,-0.55) .. (Y);
\end{tikzpicture}
\caption{}
\label{fig:19.2}
\end{figure}

图 19.2 可以表示一项随机化实验：每个时点 $k$ 的处理 $A_k$ 由研究者按依赖既往处理史与已测量协变量史的概率随机分配。若例如个体在第 $k$ 个月的处理取值这样随机分配——CD4 细胞计数低的未治疗者（$A_{k-1}=0$，$L_k=1$）以概率 0.4，CD4 细胞计数高的未治疗者（$A_{k-1}=0$，$L_k=0$）以概率 0.8，既往已治疗者不论 CD4 细胞计数如何（$A_{k-1}=1$）以概率 0.5——则我们的研究就由图 19.2 表示。在图 19.2 中，时变处理存在由已测量变量（而非不可测量变量）引起的混杂。

\begin{figure}[htbp]
\centering
\begin{tikzpicture}[causal,every node/.style={transform shape}]
  \node (L0) at (0,1.35) {$L_0$};
  \node (A0) at (2.1,1.35) {$A_0$};
  \node (L1) at (4.3,1.35) {$L_1$};
  \node (A1) at (6.5,1.35) {$A_1$};
  \node (Y)  at (8.7,1.35) {$Y$};
  \node (U0) at (0,0) {$U_0$};
  \node (U1) at (4.3,0) {$U_1$};
  \draw[ed] (U0) -- (L0);
  \draw[ed] (U0) -- (U1);
  \draw[ed] (U1) -- (L1);
  \draw[ed] (U0) -- (A0);
  \draw[ed] (L0) -- (A0);
  \draw[ed] (A0) -- (L1);
  \draw[ed] (L1) -- (A1);
  \draw[ed] (U1) -- (A1);
  \draw[ed] (A1) -- (Y);
  \draw[ed] (A0) -- (U1);
  \draw[ed] (U0) .. controls (2.2,0.10) and (3.3,0.85) .. (L1);
  \draw[ed] (L0) .. controls (1.2,2.10) and (3.1,2.10) .. (L1);
  \draw[ed] (L0) .. controls (4.6,2.95) and (5.9,2.35) .. (A1);
  \draw[ed] (A0) .. controls (3.1,2.50) and (5.5,2.50) .. (A1);
  \draw[ed] (L1) .. controls (5.4,2.00) and (7.3,2.00) .. (Y);
  \draw[ed] (A0) .. controls (4.5,3.35) and (6.8,3.35) .. (Y);
  \draw[ed] (L0) .. controls (2.9,3.85) and (6.3,3.85) .. (Y);
  \draw[ed] (U0) .. controls (2.8,-1.00) and (6.2,-1.00) .. (Y);
  \draw[ed] (U1) .. controls (5.6,-0.55) and (7.4,-0.55) .. (Y);
\end{tikzpicture}
\caption{}
\label{fig:19.3}
\end{figure}

在每个时点 $k$ 把处理随机分配给每个个体的实验，称为序贯随机化实验。因此图 19.1 与 19.2 可以表示序贯随机化实验。反之，图 19.3 不能表示随机化实验：每个时点 $k$ 的处理取值 $A_k$ 部分取决于不可测量变量 $U$，而 $U$ 是 $L_k$ 与 $Y$ 的原因；不可测量变量显然无法被研究者用来分配处理。也就是说，序贯随机化实验可以用一个含许多时点 $k=0,1,\ldots,K$、且在任何时点 $k$ 都没有从不可测量预后因子 $U$ 指向处理 $A_k$ 的直接箭头的因果图来表示。

在观察性研究中，治疗决策常取决于预后因子这类结局预测变量。因此观察性研究通常由图 19.2 或图 19.3（而非图 19.1）表示。例如，假设我们的 HIV 随访研究是一项观察性研究（不是实验），其中 CD4 细胞计数 $L_k$ 越低，病人越可能接受治疗。那么，若在每个时点 $k$，现实世界的治疗决策基于既往处理与 CD4 细胞计数史 $(\bar{A}_{k-1},\bar{L}_k)$ 的取值，而不基于任何不可测量变量 $\bar{U}_k$ 的取值，则我们的研究就由图 19.2 表示。于是，由图 19.2 表示的观察性研究与序贯随机化实验的唯一差别在于：分配概率未知（但可以从数据估计）。遗憾的是，无法凭经验判断一项观察性研究究竟是由图 19.2 还是图 19.3 的因果图表示。由图 19.3 表示的观察性研究存在不可测量混杂，这一点我们后面会讲到。

序贯随机化实验在实践中并不常用。但序贯随机化实验的概念有助于理解有效估计时变处理因果效应的一些关键条件。下一节正式给出这些条件。

\section{序贯可交换性}

如第一、第二部分所述，关于时不变处理的有效因果推断通常需要条件可交换性 $Y^a\indep A\given L$。当可交换性 $Y^a\indep A\given L$ 成立时，只要我们通过标准化、IP 加权、g 估计或其他方法对 $L$ 中的变量作适当调整，就能得到处理 $A$ 对结局 $Y$ 的因果效应的无偏估计。我们预期条件可交换性在条件随机化实验中成立——即个体接受处理的概率取决于协变量 $L$ 取值的试验。在观察性研究中，若接受处理的概率取决于已测量协变量 $L$，且在给定 $L$ 的条件下不再依赖任何不可测量的、处理与结局的共同原因，则条件可交换性成立。

类似地，时变处理的因果推断需要在每个时点调整时变协变量 $\bar{L}_k$ 以获得条件可交换性，即序贯条件可交换性。例如，在含两个时点的研究中，序贯条件可交换性是研究第一个时点与第二个时点上的条件可交换性的组合。即 $Y^g\indep A_0\given L_0$ 与 $Y^g\indep A_1\given A_0=g(L_0),L_0,L_1$。（为简洁起见，本书略去“条件”二字，径称序贯可交换性。）我们把这一组条件独立陈述称为：在任何（静态或动态）策略 $g$ 之下 $Y^g$ 的序贯可交换性，而该策略的干预涉及时变处理 $(A_0,A_1)$ 的两个分量。

序贯随机化实验——每个时点 $k$ 的处理 $A_k$ 按只依赖其既往协变量史 $\bar{L}_k$ 与处理史 $\bar{A}_{k-1}$ 的概率随机分配的实验——蕴含 $Y^g$ 的序贯可交换性。也就是说，对任何策略 $g$，在每个时点 $k$，受治疗者与未受治疗者在给定既往协变量史 $\bar{L}_k$ 以及与策略 $g$ 相容的任何观测处理史 $\bar{A}_{k-1}=g(\bar{A}_{k-2},\bar{L}_{k-1})$ 的条件下，对 $Y^g$ 是可交换的。形式上，$Y^g$ 的序贯可交换性定义为
\[
Y^g\indep A_k\given \bar{A}_{k-1}=g(\bar{A}_{k-2},\bar{L}_{k-1}),\ \bar{L}_k
\quad\text{对所有策略 } g \text{ 与 } k=0,1,\ldots,K.
\]

\mnotev{-0.6em}{对于处理史 $\bigl(\bar{A}_{k-1}=g(\bar{A}_{k-2},\bar{L}_{k-1})\bigr)$ 与其在策略 $g$ 下本会接受的处理相同（即与之相容）者，由一致性，反事实结局 $Y^g$ 等于观测结局 $Y$，因而也等于静态策略 $(a_0=\bar{A}_0,a_1=A_1)$ 下的反事实结局。}

这种形式的序贯可交换性（我们还会看到其他形式）在任何像图 19.2 那样没有从不可测量变量 $U$ 指向处理变量 $A$ 的箭头的因果图中都成立。因此，$Y^g$ 的序贯可交换性在序贯随机化实验中成立，也在这样的观察性研究中成立：其中每个时点接受治疗的概率取决于其处理史与已测量协变量史 $(\bar{A}_{k-1},\bar{L}_k)$，且在给定该历史的条件下不依赖结局的任何不可测量原因。

\mnotev{3em}{在图 19.1 中，$Y^g$ 的序贯无条件可交换性成立，即对所有静态策略 $\bar{a}$ 有 $Y^{\bar{a}}\indep A_k\given \bar{A}_{k-1}=\bar{a}_{k-1}$。无条件可交换性意味着关联即因果，即 $\E[Y^{\bar{a}}]=\E[Y\given \bar{A}=\bar{a}]$。}

也就是说，在由图 19.2 表示的观察性研究中，所有策略 $g$ 下的反事实结局均值 $\E[Y^g]$ 都可识别；而在由图 19.3 表示的观察性研究中，没有任何反事实结局均值 $\E[Y^g]$ 可识别。在由其他因果图表示的观察性研究中，某些（但不是所有）策略 $g$ 下的反事实结局均值 $\E[Y^g]$ 可识别。

例如，考虑由图 19.4 的因果图表示的观察性研究，它包含一个不可测量变量 $W_0$。在 HIV 例子中，$W_0$ 可以是时点 0 一次计划门诊就诊的指示变量，而该次就诊未记录在我们的数据库中。这种情况下，$W_0$ 会是处理 $A_0$ 与（带一定误差测量的）CD4 细胞计数 $L_1$ 的共同原因，而 $U_1$ 表示 CD4 细胞计数真实但未知的底层取值。尽管 $W_0$ 不可测量，任何静态策略 $g=\bar{a}$ 下的反事实结局均值仍可识别；但任何处理分配依赖 $L_1$ 的动态策略 $g$ 下的反事实结局均值 $\E[Y^g]$ 都不可识别。为说明为何在某些（而非所有）策略下可识别，下一节我们将使用 SWIG。

\begin{figure}[htbp]
\centering
\begin{tikzpicture}[causal,every node/.style={transform shape}]
  \node (L0) at (0,1.35) {$L_0$};
  \node (A0) at (2.1,1.35) {$A_0$};
  \node (L1) at (4.3,1.35) {$L_1$};
  \node (A1) at (6.5,1.35) {$A_1$};
  \node (Y)  at (8.7,1.35) {$Y$};
  \node (U0) at (0,0) {$U_0$};
  \node (W0) at (1.75,0.52) {$W_0$};
  \node (U1) at (4.3,0) {$U_1$};
  \draw[ed] (U0) -- (L0);
  \draw[ed] (U0) -- (U1);
  \draw[ed] (U1) -- (L1);
  \draw[ed] (L0) -- (A0);
  \draw[ed] (W0) -- (A0);
  \draw[ed] (A0) -- (L1);
  \draw[ed] (L1) -- (A1);
  \draw[ed] (A1) -- (Y);
  \draw[ed] (A0) -- (U1);
  \draw[ed] (W0) .. controls (2.6,0.70) and (3.5,1.00) .. (L1);
  \draw[ed] (U0) .. controls (2.2,0.10) and (3.3,0.85) .. (L1);
  \draw[ed] (L0) .. controls (1.2,2.10) and (3.1,2.10) .. (L1);
  \draw[ed] (L0) .. controls (4.6,2.95) and (5.9,2.35) .. (A1);
  \draw[ed] (A0) .. controls (3.1,2.50) and (5.5,2.50) .. (A1);
  \draw[ed] (L1) .. controls (5.4,2.00) and (7.3,2.00) .. (Y);
  \draw[ed] (A0) .. controls (4.5,3.35) and (6.8,3.35) .. (Y);
  \draw[ed] (L0) .. controls (2.9,3.85) and (6.3,3.85) .. (Y);
  \draw[ed] (U0) .. controls (2.8,-1.00) and (6.2,-1.00) .. (Y);
  \draw[ed] (U1) .. controls (5.6,-0.55) and (7.4,-0.55) .. (Y);
\end{tikzpicture}
\caption{}
\label{fig:19.4}
\end{figure}

除了某种形式的序贯可交换性，时变处理的因果推断还需要正值性与一致性的序贯版本。在序贯随机化实验中，序贯正值性与一致性预期都成立（见技术注记 19.2）。下面我们假定序贯正值性与一致性成立。在这三个可识别性条件下，只要我们使用对处理史与协变量史 $(\bar{A}_{k-1},\bar{L}_k)$ 作适当调整的方法——如 g 公式（标准化）、IP 加权与 g 估计——就能识别所关心策略 $g$ 下的反事实结局均值 $\E[Y^g]$。

\mnotev{-8em}{每当我们谈论因果效应的识别时，识别公式都是 g 公式（见第 21 章）。在与本书讨论无关的罕见情形中，效应可由与 g 公式相关但不相同的公式识别（例如技术注记 7.3）。}

\begin{techpoint}[动态策略的定义]
每个依赖既往处理史与协变量史的动态策略 $g=\bigl[g_0(\bar{a}_{-1},\bar{l}_0),\ldots,g_K(\bar{a}_{K-1},\bar{l}_K)\bigr]$，都对应一个只依赖既往协变量史的动态策略 $g'=\bigl[g'_0(\bar{l}_0),\ldots,g'_K(\bar{l}_K)\bigr]$。由一致性（见技术注记 19.2），个体从时点零起遵循策略 $g$ 时，其处理史、协变量史与结局史，与从时点零起遵循策略 $g'$ 时相同。特别地，$Y^g=Y^{g'}$ 且 $\bar{L}^g(K)=\bar{L}^{g'}(K)$。具体地，$g'$ 由 $g$ 递归地定义为
\[
g'_0(l_0)=g_0(\bar{a}_{-1}=0,l_0),\qquad
g'_k(\bar{l}_k)=g_k\bigl[\bar{a}_{k-1}=g'_{k-1}(\bar{l}_{k-1}),\ \bar{l}_k\bigr],
\]
其中按约定 $\bar{a}_{-1}$ 只能取零值。对任何在每个 $k$ 上处理都已不依赖既往处理史的策略 $g$，$g$ 与 $g'$ 是同一组函数。上述用 $g$ 定义 $g'$ 的方式保证：个体在观测数据中遵循策略 $g$ 直到时点 $t$（即对 $k\le t$ 有 $A_k=g_k(\bar{A}_{k-1},\bar{L}_k)$），当且仅当该个体遵循策略 $g'$ 直到 $t$（即对 $k\le t$ 有 $A_k=g'_k(\bar{L}_k)$）。
\end{techpoint}

\begin{techpoint}[时变处理的正值性与一致性]
正值性条件需要从固定版本“若 $f_L(l)\neq 0$，则对所有 $a$ 与 $l$ 有 $f_{A\given L}(a\given l)>0$”推广为序贯版本
\[
\text{若 } f_{\bar{A}_{k-1},\bar{L}_k}(\bar{a}_{k-1},\bar{l}_k)\neq 0,\quad
\text{则 } f_{A_k\given \bar{A}_{k-1},\bar{L}_k}(a_k\given \bar{a}_{k-1},\bar{l}_k)>0
\text{ 对所有 } (a_k,\bar{l}_k).
\]
在序贯随机化实验中，只要每个时点 $k$ 的随机化概率不论既往处理史与协变量史如何都不为 0 也不为 1，正值性就成立。若我们关心某个特定策略 $g$，上述正值性条件只需对与 $g$ 相容的处理史成立，即对每个 $k$，$a_k=g(\bar{a}_{k-1},l_k)$。

一致性条件也需要从固定版本“若给定个体有 $A=a$，则该个体有 $Y^a=Y$”推广为序贯版本
\[
Y^{\bar{a}}=Y^{\bar{a}^*}\ \text{若 } \bar{a}^*=\bar{a};\qquad
Y^{\bar{a}}=Y\ \text{若 } \bar{A}=\bar{a};\qquad
\bar{L}^{\bar{a}}_k=\bar{L}^{\bar{a}^*}_k=\bar{L}_k\ \text{若 } \bar{A}_{k-1}=\bar{a}_{k-1},
\]
其中 $\bar{L}^{\bar{a}}_k$ 是策略 $\bar{a}$ 下截至时点 $k$ 的反事实 $L$ 史。技术上，识别时变处理对 $Y$ 的效应只需要较弱的一致性条件：对静态策略，“若给定个体有 $\bar{A}=\bar{a}$，则该个体有 $Y^{\bar{a}}=Y$”就已足够；对动态策略，“对任何策略 $g$，若给定个体在每个时点 $k$ 都有 $A_k=g_k(\bar{A}_{k-1},\bar{L}_k)$，则 $Y^g=Y$”就已足够。不过，更强的序贯一致性是一个自然条件，我们将始终接受它。

注意，若我们预期与 $A_k=1$ 对应的干预“在第 $k$ 个月治疗”和与 $A_k=0$ 对应的干预“在第 $k$ 个月不治疗”在所有时点 $k$ 都定义得足够清楚，那么涉及 $A_k$ 的所有静态与动态策略也会同样定义清楚。
\end{techpoint}

\section{部分处理策略下的可识别性}

第 7 章给出了一个图形规则——后门准则——用以判断在特定因果图下时不变处理的可交换性是否成立。后门准则可以推广到时变处理。例如，对静态策略，处理策略的因果效应可识别的一个充分条件是：在每个时点 $k$，所有进入 $A_k$ 且不经过任何未来处理的后门路径都被阻断。

\mnotev{-0.4em}{Pearl and Robins (1995) 提出了静态策略的推广后门准则。Robins (1997a) 把该程序推广到动态策略。}

然而，推广的后门准则并不直接显示“阻断后门路径”与“序贯可交换性”之间的联系，因为该程序基于不含反事实结局的因果有向无环图。另一种检验因果效应可识别性的图形方法基于 SWIG，第 7 章也讨论过。SWIG 对时变处理尤其有用。

考虑图 19.5 与 19.6 中的因果图，它们是图 19.2 与 19.4 的简化版。我们略去了节点 $U_0$ 与 $L_0$，以及从 $A_0$ 到 $U_1$ 的箭头。此外，从 $L_1$ 到 $Y$ 的箭头不存在，因此 $L_1$ 不再是 $Y$ 的直接原因。图 19.5 与 19.6（与图 19.2、19.4 一样）的区别在于：$A_k$ 与其后（$t>k$）的协变量 $L_t$ 是否共有原因 $W_k$。

\begin{figure}[htbp]
\centering
\begin{tikzpicture}[causal,every node/.style={transform shape}]
  \node (A0) at (0,1.05) {$A_0$};
  \node (L1) at (2.3,1.05) {$L_1$};
  \node (A1) at (4.5,1.05) {$A_1$};
  \node (Y)  at (6.5,1.05) {$Y$};
  \node (U1) at (2.3,-0.5) {$U_1$};
  \draw[ed] (A0) -- (L1);
  \draw[ed] (L1) -- (A1);
  \draw[ed] (A1) -- (Y);
  \draw[ed] (U1) -- (L1);
  \draw[ed] (A0) .. controls (2.0,1.85) and (3.5,1.85) .. (A1);
  \draw[ed] (A0) .. controls (2.8,2.55) and (5.4,2.55) .. (Y);
  \draw[ed] (U1) .. controls (4.2,-0.95) and (5.9,-0.70) .. (Y);
\end{tikzpicture}
\caption{}
\label{fig:19.5}
\end{figure}

\begin{figure}[htbp]
\centering
\begin{tikzpicture}[causal,every node/.style={transform shape}]
  \node (A0) at (0,1.05) {$A_0$};
  \node (L1) at (2.3,1.05) {$L_1$};
  \node (A1) at (4.5,1.05) {$A_1$};
  \node (Y)  at (6.5,1.05) {$Y$};
  \node (W0) at (0,-0.5) {$W_0$};
  \node (U1) at (2.3,-0.5) {$U_1$};
  \draw[ed] (A0) -- (L1);
  \draw[ed] (L1) -- (A1);
  \draw[ed] (A1) -- (Y);
  \draw[ed] (W0) -- (A0);
  \draw[ed] (U1) -- (L1);
  \draw[ed] (W0) .. controls (1.0,-0.10) and (1.7,0.55) .. (L1);
  \draw[ed] (A0) .. controls (2.0,1.85) and (3.5,1.85) .. (A1);
  \draw[ed] (A0) .. controls (2.8,2.55) and (5.4,2.55) .. (Y);
  \draw[ed] (U1) .. controls (4.2,-0.95) and (5.9,-0.70) .. (Y);
\end{tikzpicture}
\caption{}
\label{fig:19.6}
\end{figure}

如本书第一部分所述，SWIG 表示特定干预下的反事实世界。图 19.7 中的 SWIG 表示：若所有人都接受静态策略 $(a_0,a_1)$（$a_0$ 与 $a_1$ 可取 0 或 1），图 19.5 的世界会是什么样。例如，图 19.7 可用来表示“始终治疗”策略（$a_0=1$，$a_1=1$）或“从不治疗”策略（$a_0=0$，$a_1=0$）下的世界。为构造这个 SWIG，我们先把处理节点 $A_0$ 与 $A_1$ 拆开。拆分后的右侧表示干预之下的处理取值，左侧表示在干预所有既往处理时会观测到的处理取值。因此 $A_0$ 的左侧恰是 $A_0$，因为没有既往处理可供干预；$A_1$ 的左侧则是把 $A_0$ 设为 $a_0$ 后会观测到的反事实处理 $A_1^{a_0}$。原因果图中所有进入某处理节点的箭头现在都指向其左侧，所有离开某处理节点的箭头现在都从右侧出发。结局变量是反事实结局 $Y^{a_0,a_1}$，协变量 $L$ 则由相应的反事实变量替换。注意我们把策略 $(a_0,a_1)$ 下与 $L_1$ 对应的反事实变量写作 $L_1^{a_0}$ 而不是 $L_1^{a_0,a_1}$，因为对未来 $A_1$ 的干预不能影响更早的 $L_1$ 的取值。

\begin{figure}[htbp]
\centering
\begin{tikzpicture}[causal,every node/.style={transform shape}]
  \node (A0) at (0,1.0) {$A_0\,|\,a_0$};
  \node (L1) at (2.4,1.0) {$L_1^{a_0}$};
  \node (A1) at (4.7,1.0) {$A_1^{a_0}\,|\,a_1$};
  \node (Y)  at (7.2,1.0) {$Y^{a_0,a_1}$};
  \node (U1) at (2.4,-0.45) {$U_1$};
  \draw[ed] (A0) -- (L1);
  \draw[ed] (L1) -- (A1);
  \draw[ed] (A1) -- (Y);
  \draw[ed] (U1) -- (L1);
  \draw[ed] (A0) to[out=58,in=125] (A1);
  \draw[ed] (A0) to[out=74,in=130] (Y);
  \draw[ed] (U1) to[out=-22,in=-105] (Y);
\end{tikzpicture}
\caption{}
\label{fig:19.7}
\end{figure}

与图 19.5 的有向无环图不同，图 19.7 的 SWIG 确实包含反事实结局，这意味着我们可以用 d 分离直观地检验可交换性。

在图 19.7 中，由 d 分离可知，对任何静态策略 $(a_0,a_1)$，$Y^{a_0,a_1}\indep A_0$ 与 $Y^{a_0,a_1}\indep A_1^{a_0}\given A_0,L_1^{a_0}$ 都成立。第二个条件独立之所以成立，是因为表面上虽存在一条开放路径 $A_1^{a_0}\leftarrow a_0\rightarrow L_1^{a_0}\leftarrow U_1\rightarrow Y^{a_0,a_1}$，但这条路径实际上被阻断了：在反事实世界中 $a_0$ 是常数，而在概率陈述中常数总是被隐含地作为条件，尽管按惯例它们不出现在条件事件中。细节见精要 19.2。

第二个条件独立 $Y^{a_0,a_1}\indep A_1^{a_0}\given A_0,L_1^{a_0}$ 按定义蕴含：在接受处理 $A_0=a_0$ 的个体子集中，$Y^{a_0,a_1}\indep A_1^{a_0}\given A_0=a_0,L_1^{a_0}$。因此由一致性可得：在图 19.5 的因果图下 $Y^{a_0,a_1}\indep A_0$ 与 $Y^{a_0,a_1}\indep A_1\given A_0=a_0,L_1$ 成立；这正对应于图 19.7 的 SWIG，而我们确实可以在那里检验可交换性。若有多个时点，我们会说
\[
Y^{\bar{a}}\indep A_k\given \bar{A}_{k-1}=\bar{a}_{k-1},\ \bar{L}_k
\quad\text{对 } k=0,1,\ldots,K.
\]

\mnotev{-0.4em}{$Y^{a_0,a_1}\indep A_1^{a_0}\given A_0=a_0,L_1^{a_0}$ 等于 $Y^{a_0,a_1}\indep A_1\given A_0=a_0,L_1$，因为由一致性，当 $A_0=a_0$ 时 $L_1^{a_0}=L_1$ 且 $A_1^{a_0}=A_1$。}

我们把上述条件称为 $Y^{\bar{a}}$ 的\emph{静态序贯可交换性}，它弱于 $Y^g$ 的序贯可交换性，因为它只要求由静态策略 $g=\bar{a}$ 索引的反事实结局 $Y^{\bar{a}}$ 与处理 $A_k$ 之间的条件独立。静态序贯可交换性足以识别任何静态策略 $g=\bar{a}$ 下的反事实结局均值。另见技术注记 19.3。

静态序贯可交换性在图 19.6 的因果图下也成立，只需对其对应的图 19.8 中的 SWIG 应用 d 分离即可验证。因此，在由图 19.6 表示的观察性研究中，我们可以识别任何静态策略 $(a_0,a_1)$ 下的反事实结局均值。

\begin{figure}[htbp]
\centering
\begin{tikzpicture}[causal,every node/.style={transform shape}]
  \node (A0) at (0,1.0) {$A_0\,|\,a_0$};
  \node (L1) at (2.4,1.0) {$L_1^{a_0}$};
  \node (A1) at (4.7,1.0) {$A_1^{a_0}\,|\,a_1$};
  \node (Y)  at (7.2,1.0) {$Y^{a_0,a_1}$};
  \node (W0) at (0,-0.45) {$W_0$};
  \node (U1) at (2.4,-0.45) {$U_1$};
  \draw[ed] (A0) -- (L1);
  \draw[ed] (L1) -- (A1);
  \draw[ed] (A1) -- (Y);
  \draw[ed] (W0) -- (A0);
  \draw[ed] (U1) -- (L1);
  \draw[ed] (W0) to[out=25,in=200] (L1);
  \draw[ed] (A0) to[out=58,in=125] (A1);
  \draw[ed] (A0) to[out=74,in=130] (Y);
  \draw[ed] (U1) to[out=-22,in=-105] (Y);
\end{tikzpicture}
\caption{}
\label{fig:19.8}
\end{figure}

让我们回到图 19.5。现在假设从 $L_1$ 到 $A_1$ 的箭头不存在。此时图 19.7 的 SWIG 中从 $L_1^{a_0}$ 到 $A_1^{a_0}$ 的箭头也会消失。于是由 d 分离可知，序贯无条件可交换性成立，因此无需 $L_1$ 的数据即可识别任何静态策略下的反事实结局均值。现在假设图 19.5 中存在从 $U_1$ 到 $A_1$ 的箭头。那么图 19.7 的 SWIG 会包含从 $U_1$ 到 $A_1^{a_0}$ 的箭头，于是任何形式的序贯可交换性都不成立，任何策略下的反事实均值都无法识别。

下面讨论动态处理策略下的 SWIG。图 19.9 表示图 19.5 的世界在动态策略 $g=[g_0,g_1(L_1)]$ 下的情形：处理 $A_0$ 被赋予固定取值 $g_0$（0 或 1），而时点 $k=1$ 的处理 $A_1$ 被赋予取值 $g_1(L_1^g)$，它取决于在时点 $k=0$ 分配处理取值 $g_0$ 之后观测到的 $L_1^g$ 的取值。例如，$g$ 可以是策略“时点 0 不治疗，仅当 CD4 细胞计数低（即 $L_1^g=1$）时在时点 1 治疗”。在该策略下，对所有人都 $g_0=0$；当 $L_1^g=1$ 时 $g_1(L_1^g)=1$，当 $L_1^g=0$ 时 $g_1(L_1^g)=0$。因此该 SWIG 包含从 $L_1^g$ 到 $g_1(L_1^g)$ 的箭头。这个箭头不是原因果图的一部分；它只存在于与该动态策略相联系的反事实世界中。因此我们把它画得与其他箭头不同，不过在评估 d 分离时仍须把它当作普通箭头对待。SWIG 中的结局是用 $L_1$ 来分配处理 $A_1$ 的动态策略 $g$ 下的反事实结局 $Y^g$。

\begin{figure}[htbp]
\centering
\begin{tikzpicture}[causal,every node/.style={transform shape}]
  \node (A0) at (0,1.0) {$A_0\,|\,g_0$};
  \node (L1) at (2.5,1.0) {$L_1^{g}$};
  \node (A1) at (5.1,1.0) {$A_1^{g}\,|\,g_1(L_1^{g})$};
  \node (Y)  at (8.1,1.0) {$Y^{g}$};
  \node (U1) at (2.5,-0.45) {$U_1$};
  \draw[ed] (A0) -- (L1);
  \draw[ed] (L1) -- (A1);
  \draw[ed] (A1) -- (Y);
  \draw[ed] (U1) -- (L1);
  \draw[->,dotted,thick] (L1) to[out=-35,in=-135] (A1);
  \draw[ed] (A0) to[out=55,in=128] (A1);
  \draw[ed] (A0) to[out=70,in=132] (Y);
  \draw[ed] (U1) to[out=-20,in=-108] (Y);
\end{tikzpicture}
\caption{}
\label{fig:19.9}
\end{figure}

对图 19.9 的 SWIG 应用 d 分离，我们发现对任何策略 $g$，$Y^g\indep A_0$ 与 $Y^g\indep A_1^g\given A_0=g_0,L_1^g$ 都成立。也就是说，$Y^g$ 的序贯可交换性成立，这意味着我们可以识别所有策略 $g$ 下的反事实结局均值（另见精要 19.2）。然而，这一结果对图 19.6 的因果图不成立。

\begin{figure}[htbp]
\centering
\begin{tikzpicture}[causal,every node/.style={transform shape}]
  \node (A0) at (0,1.0) {$A_0\,|\,g_0$};
  \node (L1) at (2.5,1.0) {$L_1^{g}$};
  \node (A1) at (5.1,1.0) {$A_1^{g}\,|\,g_1(L_1^{g})$};
  \node (Y)  at (8.1,1.0) {$Y^{g}$};
  \node (W0) at (0,-0.45) {$W_0$};
  \node (U1) at (2.5,-0.45) {$U_1$};
  \draw[ed] (A0) -- (L1);
  \draw[ed] (L1) -- (A1);
  \draw[ed] (A1) -- (Y);
  \draw[ed] (W0) -- (A0);
  \draw[ed] (U1) -- (L1);
  \draw[ed] (W0) to[out=25,in=200] (L1);
  \draw[->,dotted,thick] (L1) to[out=-35,in=-135] (A1);
  \draw[ed] (A0) to[out=55,in=128] (A1);
  \draw[ed] (A0) to[out=70,in=132] (Y);
  \draw[ed] (U1) to[out=-20,in=-108] (Y);
\end{tikzpicture}
\caption{}
\label{fig:19.10}
\end{figure}

\begin{finepoint}[SWIG 中来自干预节点的箭头]
绘制 SWIG 时，即便在单一干预世界中 $a$ 是常数、因而不能影响其他变量，我们仍画出从节点 $a$ 指向未来变量的箭头。这样做的一个理由是：它便于追踪原 DAG 中被 $A$ 直接影响的变量。第二个理由见技术注记 21.10。

为说明这一点为何重要，考虑两个因果图：1）图 7.14 的因果 DAG；2）与图 7.14 相同、但另含从 $A$ 到 $Y$ 的直接箭头的因果 DAG。如果我们不画出离开 $a$ 的箭头，这两个 DAG 对应的 SWIG 将完全相同。结果，我们就无法用 SWIG 推断出：$A$ 对 $Y$ 的因果效应在 DAG 1 中可识别（用前门公式），而在 DAG 2 中不可识别。因此，在对 SWIG 使用 d 分离时，我们需要记住：所有包含任何干预节点 $a$ 的路径都被阻断，即便我们在记号中没有显式以 $a$ 为条件。

同样的逻辑适用于确定性策略下的任何干预节点。假设我们在图 19.9 中加入基线混杂因子 $L_0$，并考虑把 $g_0$ 换成 $g_0(L_0)$ 的确定性动态策略。那么，在检验 $Y^g\indep A_1^g\given A_0,L_0,L_1^g$ 时，我们不需要另外显式以 $g_0(L_0)$ 为条件，因为给定 $L_0$ 后 $g_0(L_0)$ 成为常数。但我们需要记住：以 $L_0$ 为条件时，经过 $g_0(L_0)$ 的路径被阻断。当我们把 $A_0$ 实例化为 $g_0(L_0)$ 并利用一致性时，统计独立 $Y^g\indep A_1^g\given A_0,L_0,L_1^g$ 就变成上文描述的可交换性条件 $Y^g\indep A_1\given A_0=g_0(L_0),L_0,L_1$。

另一方面，在给每个个体从（可能依赖 $L_0$ 的）分布中分配随机处理取值 $A_0^{+,g}$ 的随机策略 $g$ 下，我们会在 SWIG 中显式加入 $A_0^{+,g}$（取代 $g_0(L_0)$），并在检验 d 分离时也把它纳入条件事件，因为 $A_0^{+,g}$ 不再由 $L_0$ 完全决定。Richardson and Robins (2013) 证明，对此类随机策略，g 公式可识别的一个必要条件是 $Y^g\indep A_1^g\given A_0,A_0^{+,g},L_0,L_1^g$ 成立。他们还考虑了依赖处理自然取值的策略（Robins et al. 2004），即在时点 $t$ 依据 $A_t^{g}$ 分配处理 $A_t^{+,g}$ 的策略，并给出了允许用扩展 g 公式识别的可交换性条件。依赖处理自然取值的策略近来被称为“修正处理政策”（modified treatment policies，Diaz et al. 2021）。
\end{finepoint}

\mnotev{-0.6em}{我们从 SWIG 读出的是 $Y^{a_0,a_1}\indep A_1^{g_0}\given A_0,L_1^{g_0}$，由一致性它蕴含 $Y^{a_0,a_1}\indep A_1\given A_0=g_0,L_1$。}

图 19.10 的 SWIG 表示图 19.6 的世界在动态处理策略 $g=[g_0,g_1(L_1)]$ 下的情形。对图 19.10 的 SWIG 应用 d 分离，我们发现 $Y^g\indep A_0$ 不成立，因为存在开放路径 $A_0\leftarrow W_0\rightarrow L_1^g\rightarrow g_1(L_1^g)\rightarrow Y^g$。也就是说，$Y^g$ 的序贯可交换性不成立，这意味着我们无法识别任何策略 $g$ 下的反事实结局均值。

总之，在由图 19.5 表示的观察性研究（或序贯随机化试验）中，$Y^g$ 的序贯可交换性成立，因此数据可用于有效估计涉及静态与动态策略的因果效应。反之，在由图 19.6 表示的观察性研究中，只有针对静态策略的较弱条件成立，因此数据可用于有效估计涉及静态策略的因果效应，但不能用于动态策略。换个角度想：在图 19.10 的 SWIG 所表示的反事实世界中，$Y^g$ 的分布取决于 $g_1(L_1^g)$ 的分布，因而取决于 $L_1^g$ 的分布；然而由于路径 $A_0\leftarrow W_0\rightarrow L_1^g$，$L_1^g$ 的分布不可识别。

最后一个例子。考虑图 19.11，它与图 19.6 相同，只多了一条从 $L_1$ 到 $Y$ 的箭头；其静态策略下对应的 SWIG 见图 19.12。用 d 分离可以证明，$Y^g$ 的序贯可交换性与 $Y^{\bar{a}}$ 的静态序贯可交换性都不成立。因此，在由图 19.11 表示的观察性研究中，我们无法用数据有效估计任何策略下的因果效应。

\begin{figure}[htbp]
\centering
\begin{tikzpicture}[causal,every node/.style={transform shape}]
  \node (A0) at (0,1.0) {$A_0$};
  \node (L1) at (2.0,1.0) {$L_1$};
  \node (A1) at (3.9,1.0) {$A_1$};
  \node (Y)  at (5.5,1.0) {$Y$};
  \node (W0) at (0,-0.45) {$W_0$};
  \node (U1) at (2.0,-0.45) {$U_1$};
  \draw[ed] (A0) -- (L1);
  \draw[ed] (L1) -- (A1);
  \draw[ed] (A1) -- (Y);
  \draw[ed] (W0) -- (A0);
  \draw[ed] (U1) -- (L1);
  \draw[ed] (W0) to[out=25,in=200] (L1);
  \draw[ed] (A0) to[out=62,in=118] (A1);
  \draw[ed] (L1) to[out=60,in=112] (Y);
  \draw[ed] (A0) to[out=78,in=126] (Y);
  \draw[ed] (U1) to[out=-24,in=-100] (Y);
\end{tikzpicture}
\caption{}
\label{fig:19.11}
\end{figure}

\begin{figure}[htbp]
\centering
\begin{tikzpicture}[causal,every node/.style={transform shape}]
  \node (A0) at (0,1.0) {$A_0\,|\,a_0$};
  \node (L1) at (2.4,1.0) {$L_1^{a_0}$};
  \node (A1) at (4.7,1.0) {$A_1^{a_0}\,|\,a_1$};
  \node (Y)  at (7.2,1.0) {$Y^{a_0,a_1}$};
  \node (W0) at (0,-0.45) {$W_0$};
  \node (U1) at (2.4,-0.45) {$U_1$};
  \draw[ed] (A0) -- (L1);
  \draw[ed] (L1) -- (A1);
  \draw[ed] (A1) -- (Y);
  \draw[ed] (W0) -- (A0);
  \draw[ed] (U1) -- (L1);
  \draw[ed] (W0) to[out=25,in=200] (L1);
  \draw[ed] (A0) to[out=58,in=125] (A1);
  \draw[ed] (L1) to[out=58,in=115] (Y);
  \draw[ed] (A0) to[out=74,in=130] (Y);
  \draw[ed] (U1) to[out=-22,in=-105] (Y);
\end{tikzpicture}
\caption{}
\label{fig:19.12}
\end{figure}

\begin{techpoint}[序贯可交换性的多种形式]
考虑一个含多个时点 $k=0,1,\ldots,K$ 的时变处理 $A_k$ 的序贯随机化实验。表示该实验的 SWIG 只是图 19.7 的加长版。可以直接从该 SWIG 读出如下条件独立：
\[
\bigl(Y^{\bar{a}},\bar{L}^{\bar{a}}_{k+1}\bigr)\indep A_k^{\bar{a}_{k-1}}\given \bar{A}^{\bar{a}_{k-2}}_{k-1},\ \bar{L}^{\bar{a}_{k-1}}_k,
\]
其中 $\bar{L}^{\bar{a}}_{k+1}$ 是从时点 $k+1$ 到随访结束的反事实协变量史。上面的条件独立
\[
\bigl(Y^{\bar{a}},\bar{L}^{\bar{a}}_{k+1}\bigr)\indep A_k^{\bar{a}_{k-1}}\given
\bar{A}^{\bar{a}_{k-2}}_{k-1}=\bar{a}_{k-1},\ \bar{L}^{\bar{a}_{k-1}}_k
\]
对特定实例 $\bar{A}^{\bar{a}_{k-2}}_{k-1}=\bar{a}_{k-1}$ 成立，其中 $\bar{a}_{k-1}$ 是策略 $\bar{a}$ 的一个分量。由一致性，最后一个条件独立陈述等于
\[
\bigl(Y^{\bar{a}},\bar{L}^{\bar{a}}_{k+1}\bigr)\indep A_k\given \bar{A}_{k-1}=\bar{a}_{k-1},\ \bar{L}_k.
\]
当这一陈述对所有 $\bar{a}$ 成立时，我们说存在序贯可交换性。有趣的是，尽管这个序贯可交换性条件只涉及静态策略 $g=\bar{a}$，它却等价于看似更强的
\[
\bigl(Y^{g},\bar{L}^{g}_{k+1}\bigr)\indep A_k\given \bar{A}_{k-1}=g(\bar{A}_{k-1},\bar{L}_k),\ \bar{L}_k
\quad\text{对所有 } g,
\]
因而（若正值性成立）足以识别任何静态与动态策略 $g$ 下的结局分布与协变量分布（Robins 1986）。这一识别来自 $(Y^{\bar{a}},\bar{L}^{\bar{a}}_{k+1})$ 与 $A_k$ 之间的联合条件独立。注意，对动态策略，序贯可交换性并不能由两个分离的独立陈述 $Y^{\bar{a}}\indep A_k\given \bar{A}_{k-1}=\bar{a}_{k-1},\bar{L}_k$ 与 $\bar{L}^{\bar{a}}_{k+1}\indep A_k\given \bar{A}_{k-1}=\bar{a}_{k-1},\bar{L}_k$ 推出。

在序贯随机化实验中，预期会有更强的条件独立成立，但它们 (i) 无法从 SWIG 读出，(ii) 对识别总体中处理策略的因果效应并非必要。例如，序贯随机化试验蕴含更强的联合独立
\[
\bigl(Y^{\bar{A}},\bar{L}^{\bar{A}}\bigr)\indep A_k\given \bar{A}_{k-1},\ \bar{L}_k,
\]
其中，对二值处理 $A_k$，$\bar{A}$ 表示全部 $2^K$ 个静态策略 $\bar{a}$ 的集合，$Y^{\bar{A}}$ 表示全部反事实结局 $Y^{\bar{a}}$ 的集合，$\bar{L}^{\bar{A}}$ 表示全部反事实协变量史的集合。借用技术注记 2.1 的术语，我们把这一联合独立条件称为\emph{完全序贯可交换性}。
\end{techpoint}

\section{时变混杂与时变混杂因子}

在观察性研究中，没有任何形式的序贯可交换性得到保证。要实现近似的可交换性，需要专家知识；专家知识会指导研究者设计研究、尽可能多地测量相关变量 $\bar{L}_k$。例如，在 HIV 研究中，专家会同意，CD4 细胞计数、病毒载量与症状这类时变变量需要被恰当测量并加以调整。

但“已测量协变量是否足以保证序贯可交换性？”这个问题永远无法确切回答。不过，我们可以用专家知识来组织关于可交换性的信念，并把它们表示在因果图中。图 19.1 至 19.4 就是概括不同情形的因果图示例。注意我们画这些因果图时没有选择（例如因失访而删失），以便在此集中讨论混杂。

考虑图 19.5。如本书第一部分那样，假设我们关心时不变处理 $A_1$ 对结局 $Y$ 的效应。我们说 $A_1$ 对 $Y$ 的效应存在混杂，因为 $A_1$ 与 $Y$ 共有原因 $U$，即 $A_1$ 与 $Y$ 之间存在经过 $U$ 的开放后门路径。要无偏地估计这一效应，我们需要调整处理 $A_1$ 对结局 $Y$ 之效应的混杂因子，如第 7 章所述。换言之，我们需要能够阻断 $A_1$ 与 $Y$ 之间所有开放的后门路径。后门路径 $A_1\leftarrow L_1\leftarrow U\rightarrow Y$ 无法通过以共同原因 $U$ 为条件来阻断，因为 $U$ 不可测量，研究者无法获得。但这条后门路径可以通过以已测量的 $L_1$ 为条件来阻断。因此，若研究者收集了所有个体的 $L_1$ 数据，则 $A_1$ 的效应就不存在不可测量混杂。于是我们说 $L_1$ 是 $A_1$ 之效应的混杂因子，尽管 $A_1$ 与 $Y$ 的实际共同原因是不可测量的 $U$（若需温习混杂与因果图，请重读 7.3 节）。

如第 7 章所述，混杂因子不必是结局的直接原因。在图 19.5 中，并不存在从混杂因子 $L_1$ 到结局 $Y$ 的箭头。此时混杂的来源（即因果混杂因子）是不可测量的共同原因 $U$。然而，由于 $L_1$ 的数据足以阻断从 $A_1$ 到 $Y$ 的后门路径、从而控制混杂，我们仍称 $L_1$ 为 $A_1$ 对 $Y$ 之效应的混杂因子。

\mnotev{-0.4em}{以对撞节点 $L_1$ 为条件后，第二条后门路径被打开：$A_1\leftarrow A_0\rightarrow L_1\leftarrow U\rightarrow Y$。只要既往处理 $A_0$ 可被研究者获得，以 $A_0$ 为条件即可安全阻断这第二条后门路径。}

现在设想一个含所有时点 $k=0,1,2,\ldots$ 的很长的因果图，其中 $L_k$ 影响其后的处理 $A_k,A_{k+1},\ldots$，并与结局 $Y$ 共有不可测量原因 $U_k$。假设我们想估计由对 $A_0,A_1,A_2$ 的干预所定义的处理策略对结局 $Y$ 的因果效应。那么在每个时点 $k$，都需要协变量史 $\bar{L}_k$ 连同处理史 $\bar{A}_{k-1}$ 来阻断处理 $A_k$ 与结局 $Y$ 之间的后门路径。因此，$\bar{A}$ 的效应无不可测量混杂，要求研究者收集所有个体的 $\bar{L}_k$ 数据。于是我们说 $\bar{L}_k$ 中的时变协变量是时变处理 $\bar{A}$ 在研究中若干（在我们的例子里是所有）时点 $k$ 上对 $Y$ 之效应的\emph{时变混杂因子}。时变混杂的更精确定义见精要 19.3。

\mnotev{-0.4em}{时变混杂因子有时也称时依混杂因子。}

遗憾的是，我们无法凭经验确认所有混杂因子（无论时不变还是时变）都已被测量。也就是说，我们无法凭经验区分无不可测量混杂的图 19.2 与有不可测量混杂的图 19.3。有趣的是，即便所有混杂因子都被正确测量与建模，多数调整方法在比较处理策略时仍可能给出有偏估计。下一章将说明为什么 g 方法是调整时变混杂因子的恰当途径。

\begin{finepoint}[时变混杂的定义]
在无选择偏倚的情况下，若 $\E[Y^{\bar{a}}]\neq \E[Y\given \bar{A}=\bar{a}]$，即若（与事实相反地）研究中的所有个体都遵循策略 $\bar{a}$ 时的平均结局，不同于实际研究中遵循策略 $\bar{a}$ 的那部分个体的平均结局，我们就说涉及 $\E[Y^{\bar{a}}]$ 的因果效应存在混杂。

若 $\E[Y^{\bar{a}}\given L_0]=\E[Y\given \bar{A}=\bar{a},L_0]$，我们就说混杂完全是时不变的（即完全可归因于基线协变量）；当图 19.2 中指向 $A_1$ 的箭头只来自 $A_0$ 与 $L_0$ 时，就是这种情形。反之，若可识别性条件成立，但 $\E[Y^{\bar{a}}\given L_0]\neq \E[Y\given \bar{A}=\bar{a},L_0]$，我们就说存在时变混杂。若可识别性条件不成立，如图 19.3，我们就说存在不可测量混杂。

无时变混杂的一个充分条件，是 $Y^{\bar{a}}$ 的无条件序贯可交换性，即 $Y^{\bar{a}}\indep A_k\given \bar{A}_{k-1}=\bar{a}_{k-1}$。该条件在序贯随机化实验中成立，例如图 19.1 所表示的那类实验：其中每个时点 $k$ 的处理 $A_k$ 按只依赖既往处理史 $\bar{A}_{k-1}$ 取值的概率随机分配。事实上，图 19.1 的因果图可以大大简化。为此，首先注意 $L_1$ 不是图中任何两个节点的共同原因，故可从图中略去。$L_1$ 一旦去掉，$L_0$ 与 $U_1$ 也都可以略去，因为它们不再是图中两个节点的共同原因。在去掉 $L_0$、$L_1$ 与 $U_1$ 的图中，节点 $U_0$ 也可略去。也就是说，图 19.1 的因果图可以简化为只含节点 $A_0$、$A_1$ 与 $Y$。
\end{finepoint}
